% =====================================================================
%  CARNET D'ENTRAINEMENT — FICHE 11 — THEME 2
%  Tuto : Oxydant, réducteur et demi-équations
% =====================================================================
\documentclass[11pt,a4paper]{article}
\usepackage[utf8]{inputenc}
\usepackage[T1]{fontenc}
\usepackage{lmodern}
\usepackage[french]{babel}
\usepackage{amsmath,amssymb}
\usepackage[version=4]{mhchem}
\usepackage{siunitx}
\sisetup{output-decimal-marker={,},per-mode=symbol}
\usepackage{xcolor}
\usepackage{array}
\usepackage{booktabs}
\usepackage{enumitem}
\usepackage[most]{tcolorbox}
\usepackage{tikz}
\usetikzlibrary{arrows.meta,positioning}
\usepackage[a4paper,margin=2.1cm,headheight=26pt,headsep=10pt]{geometry}
\usepackage{fancyhdr}
\usepackage{titlesec}
\usepackage[hidelinks]{hyperref}

\definecolor{primary}{RGB}{146,95,160}
\definecolor{primarydark}{RGB}{92,52,104}
\definecolor{defcol}{RGB}{0,114,178}
\definecolor{thmcol}{RGB}{0,140,120}
\definecolor{excol}{RGB}{0,135,90}
\definecolor{attncol}{RGB}{200,40,55}
\definecolor{methcol}{RGB}{90,90,100}
\definecolor{oxcol}{RGB}{200,40,55}
\definecolor{redcol}{RGB}{0,110,180}

\newcommand{\NO}{\ensuremath{\mathrm{N.O.}}}

\tcbset{boxbase/.style={enhanced,breakable,boxrule=0.9pt,arc=2.5pt,
   left=3mm,right=3mm,top=2.3mm,bottom=2.3mm,fonttitle=\bfseries,coltitle=white,
   toptitle=0.6mm,bottomtitle=0.6mm}}
\newtcolorbox{defi}[1][]{boxbase,colback=defcol!5,colframe=defcol,title={\textsc{Définition}\ifx\relax#1\relax\relax\else\textmd{ --- \itshape #1}\fi}}
\newtcolorbox{regle}[1][]{boxbase,colback=thmcol!6,colframe=thmcol,title={\textsc{Règle}\ifx\relax#1\relax\relax\else\textmd{ --- \itshape #1}\fi}}
\newtcolorbox{exemple}[1][]{boxbase,colback=excol!5,colframe=excol,title={\textsc{Exemple}\ifx\relax#1\relax\relax\else\textmd{ : \itshape #1}\fi}}
\newtcolorbox{attention}[1][]{boxbase,colback=attncol!6,colframe=attncol,title={\textsc{Attention}\ifx\relax#1\relax\relax\else\textmd{ : \itshape #1}\fi}}
\newtcolorbox{methode}[1][]{boxbase,colback=methcol!7,colframe=methcol,sharp corners,title={\textsc{Méthode}\ifx\relax#1\relax\relax\else\textmd{ : \itshape #1}\fi}}

\pagestyle{fancy}\fancyhf{}
\fancyhead[L]{\small\textcolor{primarydark}{\textbf{Fiche 11} — Tuto Thème 2}}
\fancyhead[R]{\small\textcolor{primarydark}{\textsc{Oxydant, réducteur, demi-équations}}}
\fancyfoot[C]{\small\thepage}
\renewcommand{\headrulewidth}{0.8pt}
\renewcommand{\headrule}{\hbox to\headwidth{\color{primary}\leaders\hrule height \headrulewidth\hfill}}
\titleformat{\section}{\large\bfseries\color{primarydark}}{\thesection}{0.6em}{}

\begin{document}
\begin{center}
{\LARGE\bfseries\color{primarydark}Thème 2 — Oxydant, réducteur et demi-équations}\\[3pt]
{\color{primary}\rule{9cm}{1.2pt}}\\[4pt]
{\small\itshape À lire avant les exercices — lire une réaction redox comme un transfert d'électrons.}
\end{center}

\vspace{-2pt}
\section*{1. Oxydation, réduction : une histoire de \NO}
\begin{defi}[oxydation et réduction]
\begin{itemize}[leftmargin=1.5em,topsep=1pt,itemsep=1.5pt]
  \item Un élément est \textbf{oxydé} si son \NO\ \textbf{augmente} : il \textbf{perd} des électrons.
  \item Un élément est \textbf{réduit} si son \NO\ \textbf{diminue} : il \textbf{capte} des électrons.
\end{itemize}
Les deux phénomènes sont \textbf{toujours simultanés} : si une espèce s'oxyde, une autre se réduit.
\end{defi}

\begin{defi}[oxydant et réducteur]
L'\textbf{\textcolor{oxcol}{oxydant}} est l'espèce qui \textbf{capte} les électrons : c'est celle qui
est \textbf{réduite}. \\
Le \textbf{\textcolor{redcol}{réducteur}} est l'espèce qui \textbf{cède} les électrons : c'est celle
qui est \textbf{oxydée}.
\end{defi}

\begin{center}
\small\framebox{\parbox{0.92\textwidth}{\centering
Moyen mnémotechnique : \og{}\textbf{on oxyde un réducteur, on réduit un oxydant}\fg{}. \\
Un \textbf{réducteur} d\textbf{onne} des électrons ; un \textbf{oxydant} en \textbf{capte}.}}
\end{center}

\section*{2. Le couple Ox/Red et la demi-équation}
\begin{defi}[couple oxydo-réducteur]
Un couple $\mathrm{Ox}/\mathrm{Red}$ relie une forme oxydée et une forme réduite par une
\textbf{demi-équation électronique} :
\[
\mathrm{Ox} + n\,\ce{e^-} \;\rightleftharpoons\; \mathrm{Red}.
\]
\end{defi}

\begin{attention}[où placer les électrons ?]
Les électrons se placent \textbf{toujours du côté de l'oxydant} : c'est lui qui les capte pour
devenir le réducteur. Ex. couple \ce{Fe^{3+}}/\ce{Fe^{2+}} : $\ce{Fe^{3+} + e^- <=> Fe^{2+}}$
(et surtout pas l'inverse).
\end{attention}

\begin{regle}[nombre d'électrons échangés]
\[
n(\ce{e^-}) = \Delta\NO = \NO(\text{oxydant}) - \NO(\text{réducteur}).
\]
Dans un couple, l'\textbf{oxydant} est la forme au \NO\ le \emph{plus élevé}, le \textbf{réducteur}
celle au \NO\ le \emph{plus bas}.
\end{regle}

\begin{exemple}[le couple \ce{NO3^-}/\ce{NO}]
$\NO(\ce{N})$ passe de $+\mathrm{V}$ (dans \ce{NO3^-}) à $+\mathrm{II}$ (dans \ce{NO}) : différence
$n=5-2=3$ électrons. Comme \ce{NO3^-} est le plus oxydé, il est l'\textbf{oxydant}, et l'électron va
de son côté : $\ce{NO3^- + 3 e^- <=> NO}$.
\end{exemple}

\section*{3. Reconnaître une réaction redox}
\begin{methode}[une réaction est-elle redox ?]
Une réaction est \textbf{redox} si (et seulement si) \textbf{au moins un élément change de \NO}
entre réactifs et produits. Sinon (acide-base, précipitation, complexation\ldots), ce n'est
\textbf{pas} du redox.
\end{methode}

\begin{exemple}[trier des réactions]
\begin{itemize}[leftmargin=1.5em,topsep=1pt,itemsep=1pt]
  \item $\ce{NaCl + AgNO3 -> AgCl + NaNO3}$ : aucun \NO\ ne change (précipitation) $\Rightarrow$ \textbf{pas redox}.
  \item $\ce{Cu + 2 Ag+ -> Cu^{2+} + 2 Ag}$ : Cu ($0\to+\mathrm{II}$), Ag ($+\mathrm{I}\to 0$) $\Rightarrow$ \textbf{redox}.
\end{itemize}
\end{exemple}

\section*{4. Combiner deux demi-équations}
Une oxydation et une réduction n'existent jamais seules. Pour obtenir l'\textbf{équation globale},
on multiplie chaque demi-équation pour que le \textbf{nombre d'électrons soit identique} des deux
côtés, puis on additionne : les électrons s'éliminent.

\begin{exemple}[nombres d'électrons différents]
Couple 1 : $\ce{A -> A^{2+} + 2 e^-}$ (2 e$^-$). Couple 2 : $\ce{B + 3 e^- -> B^{3-}}$ (3 e$^-$).
On multiplie par 3 et par 2 (soit 6 e$^-$ chacun) :
\[
\ce{3 A + 2 B -> 3 A^{2+} + 2 B^{3-}}.
\]
\end{exemple}

\begin{attention}[ions spectateurs]
Dans une équation en solution, les ions qui ne changent \emph{pas} de \NO\ (souvent \ce{Na+},
\ce{K+}, \ce{Cl^-}\ldots) sont des \textbf{ions spectateurs} : ils n'apparaissent pas dans
l'équation redox simplifiée.
\end{attention}
\end{document}
